\subsection{Начало работы с WebGPU}

Поскольку WebGPU является относительно новой технологией, для начала работы с ней потребуется выполнить несколько шагов по настройке \cite{webgpufundamentals}.

Во-первых, необходимо убедиться, что браузер поддерживает WebGPU. На момент написания этой работы, поддержка WebGPU была доступна в последних версиях браузеров Chromium при включении экспериментальных функций \cite{webgpustatus}.

Во-вторых, необходимо настроить графический конвеер для выполнения вычислений на GPU. В WebGPU это включает создание устройства GPU, выделение буферов для хранения данных и написание вычислительных шейдеров на языке WGSL (WebGPU Shading Language).

В-третьих, потребуется интегрировать динамические параметры симуляции, такие как скорость течения, вязкость и проч., а также эмулировать внешние силы, чтобы обеспечить интерактивную и наглядную визуализацию течения жидкости.

В результате этих шагов будет создана основа для реализации численного алгоритма моделирования жидкости с использованием WebGPU. Примерный план выполнения программы данных будет выглядеть следующим образом:
\begin{enumerate}
    \item Инициализация WebGPU
          \begin{enumerate}
              \item Создание и настройка модулей WebGPU.
              \item Создание буферов для хранения данных симуляции на GPU.
              \item Подключение шейдеров для выполнения вычислений.
              \item Создания графического конвейера для визуализации результатов.
              \item Запуск цикла рендеринга и обновления симуляции.
          \end{enumerate}
    \item Основной цикл симуляции
          \begin{enumerate}
              \item Обновление входных данных симуляции на CPU.
              \item Передача этих данных на GPU через буферы WebGPU.
              \item Выполнение вычислительных шейдеров, обновления состояния симуляции на GPU.
          \end{enumerate}
\end{enumerate}



\subsection{Анализ алгоритма и выбор шейдеров}
Как упоминалось уже ранее, метод Stable Fluids основан на приближении \eqref{eq:NS_projection} уравнения Навье-Стокса. Выбирая $\Delta t$, мы можем построить симуляцию с помощью разложения на четыре этапа: добавление силы, адвекция, диффузия и проекция. Каждый из этих этапов может быть реализован с помощью отдельного шейдера, что позволяет просто и эффективно использовать вычисления на GPU.

В классическом варианте метода Stable Fluids\cite{stamstablefluids} каждый этап можно записать следующим образом:
\begin{equation}\label{eq:stages_pipeline}
    \mathbf{w}_0(\mathbf{x})
    \xrightarrow{\text{сила}}
    \mathbf{w}_1(\mathbf{x})
    \xrightarrow{\text{адвекция}}
    \mathbf{w}_2(\mathbf{x})
    \xrightarrow{\text{диффузия}}
    \mathbf{w}_3(\mathbf{x})
    \xrightarrow{\text{проекция}}
    \mathbf{w}_4(\mathbf{x})
\end{equation}


\[
    \begin{aligned}
        \mathbf{w}_1(\mathbf{x}) & = \mathbf{w}_0(\mathbf{x}) + \Delta t\,\mathbf{f}(\mathbf{x}), \\
        \mathbf{w}_2(\mathbf{x}) & = \operatorname{Advect}(\mathbf{w}_1,\Delta t),                \\
        \mathbf{w}_3(\mathbf{x}) & = \operatorname{Diffuse}(\mathbf{w}_2,\nu,\Delta t),           \\
        \mathbf{w}_4(\mathbf{x}) & = \mathcal{P}\bigl(\mathbf{w}_3(\mathbf{x})\bigr),
    \end{aligned}
\]

где каждый из операторов может быть реализован с помощью отдельного шейдера. Например, адвекция может быть реализована с помощью обратного трассирования частиц, диффузия \textemdash{} с помощью метода Якоби, а проекция \textemdash{} с помощью решения уравнения Пуассона.

Также весьма полезно будет реализовать дополнительные шейдеры для визуализации объектов симуляции, а также результатов симуляций такие как отрисовка векторного поля скорости или визуализация давления.

